% Bagian: sets-functions-relations
% Bab: size-of-sets
% Subbagian: zig-zag

\documentclass[../../../include/open-logic-section]{subfiles}

\begin{document}

\olfileid[id]{sfr}{siz}{zigzag}

\olsection{Metode Zig-Zag Cantor}

\begin{explain}
Kita telah meninjau beberapa enumerasi yang ``mudah''. Sekarang kita akan
meninjau sesuatu yang sedikit lebih sulit. Perhatikan himpunan pasangan
bilangan asli\oliflabeldef{sfr:set:pai:sec}{, yang telah kita definisikan dalam
\olref[set][pai]{sec} sebagai berikut:}{yang didefinisikan oleh:}
\[
\Nat \times \Nat = \Setabs{\tuple{n,m}}{n,m \in \Nat}
\]
Kita dapat menyusun pasangan-pasangan terurut ini dalam sebuah \emph{larik},
sebagai berikut:
\[
\begin{array}{ c | c | c | c | c | c}
& \mathbf 0 & \mathbf 1 & \mathbf 2 & \mathbf 3 & \dots \\
\hline
\mathbf 0 & \tuple{0,0} & \tuple{0,1} & \tuple{0,2} & \tuple{0,3} & \dots \\
\hline
\mathbf 1 & \tuple{1,0} & \tuple{1,1} & \tuple{1,2} & \tuple{1,3} & \dots \\
\hline
\mathbf 2 & \tuple{2,0} & \tuple{2,1} & \tuple{2,2} & \tuple{2,3} & \dots \\
\hline
\mathbf 3 & \tuple{3,0} & \tuple{3,1} & \tuple{3,2} & \tuple{3,3} & \dots \\
\hline
\vdots & \vdots & \vdots & \vdots & \vdots & \ddots\\
\end{array}
\]
Jelas bahwa setiap pasangan terurut dalam $\Nat \times \Nat$ akan muncul
tepat satu kali dalam larik tersebut. Secara khusus, $\tuple{n,m}$ akan muncul
pada baris ke-$n$ dan kolom ke-$m$. Namun, bagaimana kita menyusun
anggota-anggota suatu larik semacam itu menjadi daftar ``satu dimensi''? Pola
dalam larik berikut menunjukkan salah satu caranya (meskipun tentu saja ada
banyak pilihan lain):
\[
\begin{array}{ c | c | c | c | c | c | c}
& \mathbf 0 & \mathbf 1 & \mathbf 2 & \mathbf 3 & \mathbf 4 &\dots \\
\hline
\mathbf 0 & 0  & 1& 3 & 6& 10 &\ldots \\
\hline
\mathbf 1 &2 & 4& 7 & 11 & \dots &\ldots \\
\hline
\mathbf 2 & 5 & 8 & 12 & \ldots & \dots&\ldots \\
\hline
\mathbf 3 & 9 & 13 & \ldots & \ldots & \dots & \ldots \\
\hline
\mathbf 4 & 14 & \ldots & \ldots & \ldots & \dots & \ldots \\
\hline
\vdots & \vdots & \vdots & \vdots & \vdots&\ldots & \ddots\\
\end{array}
\]\noindent
Pola ini disebut \emph{metode zig-zag Cantor}. Pola tersebut mengenumerasi
$\Nat \times \Nat$ sebagai berikut:
\[
\tuple{0,0}, \tuple{0,1}, \tuple{1,0}, \tuple{0,2}, \tuple{1,1},
\tuple{2,0}, \tuple{0,3}, \tuple{1,2}, \tuple{2,1}, \tuple{3,0}, \dots
\]
Dengan demikian, diperoleh hasil berikut:
\end{explain}

\begin{prop}\ollabel{natsquaredenumerable}
$\Nat \times \Nat$ adalah !!{enumerable}.
\end{prop}

\begin{proof}
Misalkan $f \colon \Nat \to \Nat\times\Nat$ memetakan setiap $k \in \Nat$
ke tupel $\tuple{n,m} \in \Nat \times \Nat$ sedemikian sehingga $k$ adalah
nilai pada baris ke-$n$ dan kolom ke-$m$ dalam larik zig-zag Cantor.
\end{proof}

\begin{explain}
Teknik ini juga dapat diperumum dengan cukup baik. Sebagai contoh, kita dapat
menggunakannya untuk mengenumerasi himpunan tripel terurut bilangan asli,
yakni:
\[
\Nat \times \Nat \times \Nat = \Setabs{\tuple{n,m,k}}{n,m,k \in \Nat}
\]
Kita memandang $\Nat \times \Nat \times \Nat$ sebagai hasil kali Kartesius
dari $\Nat \times \Nat$ dengan $\Nat$, yaitu,
\[
\Nat^3 = (\Nat \times \Nat) \times \Nat =
\Setabs{\tuple{\tuple{n,m},k}}{n, m, k
  \in \Nat }
\]
dan karena itu kita dapat mengenumerasi $\Nat^3$ dengan sebuah larik: satu
sumbunya diberi label menurut enumerasi $\Nat$, sedangkan sumbu lainnya
menurut enumerasi $\Nat^2$:
\[
\begin{array}{ c | c | c | c | c | c}
& \mathbf 0 & \mathbf 1 & \mathbf 2 & \mathbf 3 & \dots \\
\hline
\mathbf{\tuple{0,0}} & \tuple{0,0,0} & \tuple{0,0,1} & \tuple{0,0,2} & \tuple{0,0,3} & \dots \\
\hline
\mathbf{\tuple{0,1}} & \tuple{0,1,0} & \tuple{0,1,1} & \tuple{0,1,2} & \tuple{0,1,3} & \dots \\
\hline
\mathbf{\tuple{1,0}} & \tuple{1,0,0} & \tuple{1,0,1} & \tuple{1,0,2} & \tuple{1,0,3} & \dots \\
\hline
\mathbf{\tuple{0,2}} & \tuple{0,2,0} & \tuple{0,2,1} & \tuple{0,2,2} & \tuple{0,2,3} & \dots\\
\hline
\vdots & \vdots & \vdots & \vdots & \vdots & \ddots \\
\end{array}
\]
Jadi, dengan menggunakan metode seperti metode zig-zag Cantor, kita juga
dapat memperoleh enumerasi~$\Nat^3$. Proses ini dapat diteruskan sehingga
kita memperoleh enumerasi $\Nat^n$ untuk setiap bilangan asli $n$. Dengan
demikian, kita mempunyai:
\end{explain}

\begin{prop}
$\Nat^n$ adalah !!{enumerable}, untuk setiap $n \in \Nat$.
\end{prop}

\begin{prob}\label{sfr:siz:zigzag:prob:posint-n}
Tunjukkan bahwa $(\PosInt)^n$ adalah !!{enumerable}, untuk setiap $n \in \Nat$.
\end{prob}

\begin{prob}\label{sfr:siz:zigzag:prob:posint-star}
  Tunjukkan bahwa $(\PosInt)^*$ adalah !!{enumerable}. Anda boleh menggunakan
  \cref{sfr:siz:zigzag:prob:posint-n}.
\end{prob}

\end{document}
